% ============================================================
%  Summer Report — Prediction-Based Lossy Compression
%  (DEIM / Gaussian Process methods on ISABEL Hurricane & SST)
%
%  HOW TO USE THIS FILE
%  ---------------------
%  Every section below contains an OUTLINE block in red italic text
%  (\prompt{...}) instead of drafted prose. These are prompts/reminders
%  of what to cover, written by working through the code and past
%  discussion with Claude. Write your own prose to replace/surround
%  them, then delete the \prompt{...} lines before final submission.
%
%  To use in Overleaf:
%    1. New Project -> Blank Project in Overleaf
%    2. Upload this file, replacing main.tex (or set this as the
%       main document under Menu -> Settings)
%    3. Upload figures referenced below into the same project
%       (see the "FIGURES NEEDED" notes in each section)
% ============================================================
\documentclass[11pt]{article}

\usepackage[margin=1in]{geometry}
\usepackage{graphicx}
\usepackage{amsmath, amssymb}
\usepackage{booktabs}
\usepackage{hyperref}
\usepackage{xcolor}
\usepackage{caption}

% Visible outline/prompt helper — delete \prompt{} calls before final submission.
\newcommand{\prompt}[1]{{\color{red}\itshape [PROMPT: #1]}\par\medskip}

\title{Prediction-Based Lossy Compression for Scientific Data:\\
       Gaussian Process and DEIM Methods on the ISABEL Hurricane and NOAA SST Datasets}
\author{Jessie Chen \\ Argonne National Laboratory}
\date{Summer 2026}

\begin{document}
\maketitle

% ============================================================
\begin{abstract}
\prompt{2--4 sentences. State the overall goal (using prediction models
as a preprocessing/predictor step ahead of quantization, to beat
general-purpose compressors like SZ2/SZ3/ZFP on structured scientific
fields), name the two datasets (ISABEL hurricane, NOAA OI SST), and
give the one-sentence status of where things currently stand
(some methods beat the baseline, the newest method does not yet --
be explicit that this is stated plainly, not hidden).}
\end{abstract}

% ============================================================
\section{Motivation and Background}
\label{sec:motivation}

\prompt{Explain why this project exists. Suggested points to hit:
\begin{itemize}
  \item What SZ2/ZFP/SZ3 do at a high level (predict each point from
        its neighbors via a local predictor, e.g. Lorenzo; quantize the
        residual; entropy-code; fall back to raw storage when the
        predictor fails on a point).
  \item Why a \emph{global} predictor (GP regression or a low-rank
        SVD/DEIM basis, fit from a small number of sensor readings)
        might do better than a \emph{local} predictor like Lorenzo on
        fields with strong global spatial structure (e.g. a hurricane's
        vortex, or SST's large-scale ocean currents).
  \item What "the compressor" means in this project's framework:
        predict the full field from $k$ sensor readings $\to$ residual
        $=$ true $-$ predicted $\to$ quantize/entropy-code the residual
        the same way SZ does, so comparisons are apples-to-apples.
  \item One or two sentences citing the two real datasets used
        (ISABEL Hurricane Uf48, 100$\times$500$\times$500 wind field;
        NOAA OI SST V2 weekly means, 1727 snapshots).
\end{itemize}}

% ============================================================
\section{Most Recent Method: Progressive GP / Batch-Sequential Anchor Placement}
\label{sec:progressive-gp}

\prompt{This is the section you're actively improving --- write it last,
once the outlier-fallback fix is implemented and re-run. Keep the
structure below either way, since you agreed to report this
transparently.}

\subsection{Idea and Algorithm}
\prompt{Describe the round-based idea: each round, place $k$ new
anchor/sensor points via randomized pivoted Cholesky (targeting
high-uncertainty regions), fit a GP on all accumulated sensors, predict
everywhere, and mark points within an acceptance band as done. Mention
this is implemented in \texttt{hurricane\_progressive\_gp.py}.
FIGURE NEEDED: algorithm diagram / the "The Idea" and "Algorithm"
slides you already built in \texttt{slides\_anchor\_point.tex} can be
reused or adapted here.}

\subsection{Current Status: Comparison Against SZ3}
\prompt{State plainly, with numbers, where this stands as of the most
recent run (from \texttt{sz3\_gp\_comparison.py} / \texttt{sz3\_gp\_comparison.png}):
SZ3 reaches PSNR 83.4 dB with 100\% of voxels within the 0.01 m/s error
bound; Progressive GP reaches PSNR 36.6 dB with only 28.2\% of voxels
within bound after 10 rounds (2,500 sensors, $\sim$0.1\% of the domain).
Note the trend is favorable round-over-round (entropy falling,
PSNR rising monotonically) even though the absolute gap is large.
FIGURE NEEDED: \texttt{sz3\_gp\_comparison.png} (4-panel comparison),
plus \texttt{frames\_hists/hist\_010.png} or similar for the error
histogram.}

\subsection{Diagnosis: Why the Gap}
\prompt{Explain the root cause you identified: the current algorithm
has no fallback/outlier mechanism. SZ2/SZ3 guarantee their error bound
by storing any point their predictor can't hit raw/exactly; Progressive
GP currently just leaves unresolved points unresolved after the round
budget runs out. Contrast this against the histogram plot's own
visualization issue you caught (percentile-clipped x-axis in
\texttt{\_compute\_edges()}, and the accept band being visually
imperceptible against the plotted range) -- note that finding as a
secondary, presentation-only bug, separate from the real algorithmic gap.}

\subsection{Fix and (Once Complete) Updated Results}
\prompt{Describe the SZ-style quantize-with-outlier-fallback fix:
quantize every point's residual against \texttt{EB} directly (not the
looser \texttt{ACCEPT\_BOUND} used mid-loop); anything that would clip
gets stored raw instead. This guarantees 100\% coverage by construction.
Once you've re-run \texttt{sz3\_gp\_comparison.py} with this fix, replace
this prompt with the updated numbers/figure and be honest about the new
bit-rate cost of the outlier fraction.}

% ============================================================
\section{Earlier Methods: Baseline-Compared Reconstruction Approaches}
\label{sec:earlier-methods}

\prompt{One-paragraph transition: these are the methods that already
have completed baseline comparisons (vs. SZ2/ZFP), built up over the
summer before the progressive-GP idea. State the throughline: Kriging
and DEIM are the two core reconstruction methods; they were packaged as
libpressio-compatible external compressors; Kriging was extended to 3D;
and three head-to-head rate-distortion scripts tie everything together.}

\subsection{Kriging / Gaussian Process Reconstruction}
\prompt{From \texttt{kriging\_hurricane.py}: Simple Kriging with an
RBF/Mat\'ern kernel, sensors placed via GKS (randomly-pivoted or
greedily-pivoted Cholesky), reconstructing held-out ISABEL levels from a
sparse sensor set. Mention this is a standalone reconstruction-quality
demo (no baseline comparison itself) -- the baseline comparison happens
later in \texttt{rate\_distortion\_comparison.py} et al. Include your own
representative reconstruction numbers/figure if you have one saved
(e.g. \texttt{kriging\_results.png}).}

\subsection{DEIM Reconstruction}
\prompt{From \texttt{deim\_hurricane\_v2.py}: Q-DEIM (SVD/POD basis +
QR-pivoted sensor placement), sensor count chosen automatically from a
POD-energy error tolerance. Same note as above -- standalone
reconstruction demo, baseline comparison comes later.}

\subsection{Packaging as libpressio Compressors}
\prompt{From \texttt{lp\_kriging\_compressor.py} and
\texttt{lp\_deim\_compressor.py}: train/compress/decompress CLI wrappers
around the two methods above, registered via libpressio's
\texttt{external:command} protocol so they can be driven identically to
SZ/ZFP by any libpressio-based tool. Note this is an interoperability
packaging, not something the other scripts import directly (they
reimplement the methods inline instead).}

\subsection{3D Extension: Kronecker GP}
\prompt{From \texttt{hurricane\_rd\_3d.py}: extends Kriging to a
separable 3D kernel ($K_{\text{spatial}} \otimes K_{\text{vertical}}$)
so a single model reconstructs the full volume at once, compared
directly against SZ2/ZFP on the TC (temperature) field. State the
result numbers/figure from \texttt{poster\_rd\_psnr\_3D\_Hurricane.png}.}

% ============================================================
\section{Rate-Distortion Comparisons Against SZ2/ZFP}
\label{sec:rd-comparisons}

\prompt{This section is where the actual baseline numbers live --
you said "we have everything we need" for this part, so this is
likely the most straightforward section to write first.}

\subsection{ISABEL Hurricane (CLOUD/QVAPOR fields)}
\prompt{From \texttt{rate\_distortion\_comparison.py}: SZ2 vs ZFP vs
DEIM-2D vs Kriging-2D (T-DEIM and MultiGP results exist here too but
are excluded per Section~\ref{sec:excluded}). Report where DEIM-2D and
Kriging-2D land relative to SZ2/ZFP on the PSNR-vs-compression-ratio
curve -- do they win at any operating point? Cite specific CR/PSNR
numbers from \texttt{rd\_results\_\{FIELD\_TAG\}.csv}.
FIGURES NEEDED: \texttt{rd\_psnr\_cr\_*.png}, \texttt{poster\_rd\_psnr\_*.png}.}

\subsection{ISABEL Hurricane (TC / temperature field)}
\prompt{From \texttt{hurricane\_rd\_comparison.py}: same four-method
comparison (SZ2, ZFP, DEIM-2D, Kriging-2D) but on the TC field with no
land mask. Note this is where the 3D Kronecker result from
Section~\ref{sec:earlier-methods} extends the comparison further.}

\subsection{NOAA OI SST}
\prompt{From \texttt{sst\_rd\_comparison.py}: SZ2/SZ3/ZFP vs DEIM-2D vs
Kriging-2D on 1727 weekly SST snapshots (ocean pixels only). This
dataset has by far the most residual-coding variants tested
(\texttt{-L2}, \texttt{-LP}, \texttt{-SO}, \texttt{-TT}, hybrid
\texttt{+SZ2}/\texttt{+ZFP}) -- decide how much of that detail belongs
in the main report vs. an appendix.}

\subsection{Supporting Diagnostics}
\prompt{Briefly mention \texttt{lorenzo\_vs\_global\_diagnostic.py} and
\texttt{residual\_autocorr\_diagnostic.py} if their findings support your
main argument (e.g. global predictors leaving less spatially-correlated
residual than Lorenzo) -- these were diagnostic-only, no new compression
run, so a sentence or two is probably enough rather than a full subsection.}

% ============================================================
\section{Discussion}
\label{sec:discussion}

\prompt{Synthesize Sections~\ref{sec:earlier-methods}--\ref{sec:rd-comparisons}:
where do the global-predictor methods (DEIM/Kriging) actually beat
SZ2/ZFP, and where do they not? What's the mechanism (e.g. GP/DEIM wins
at low bit rates / high compression ratios where a good global
prediction saves more than local prediction; loses at high bit rates
where SZ's cheap fast local predictor is good enough already)? Tie this
back to why the progressive-GP idea in Section~\ref{sec:progressive-gp}
was worth trying (targeting the sensor-budget inefficiency of the
static methods) even though it isn't there yet.}

% ============================================================
\section{Excluded From This Report: T-DEIM and MultiGP}
\label{sec:excluded}

\prompt{One paragraph, as agreed. State that \texttt{tdeim\_hurricane.py}
/ \texttt{lp\_tdeim\_compressor.py} and \texttt{multigp\_hurricane.py} /
\texttt{lp\_multigp\_compressor.py} were developed and produce
reconstructions, but have not yet been run through the same SZ2/ZFP
baseline comparison pipeline as the methods above, so they're left out
of the quantitative discussion here. Note briefly what each does
(T-DEIM: sensors placed anywhere in the full 3D volume rather than a
single slice; MultiGP: joint multi-field reconstruction via a Linear
Model of Coregionalization) and flag them as candidates for a follow-up
comparison.}

% ============================================================
\section{Conclusion and Next Steps}
\label{sec:conclusion}

\prompt{Summarize: which methods currently beat the SZ2/ZFP/SZ3
baseline and by how much; what the progressive-GP method still needs
(outlier fallback, then a fair rerun); what's still open for T-DEIM/MultiGP.
List concrete next steps -- e.g. finish the quantize+outlier fix, rerun
the SZ3 comparison, decide whether the sensor-budget/entropy trend from
Section~\ref{sec:progressive-gp} is compelling enough to pursue further
even after accounting for the guaranteed-bound cost.}

% ============================================================
% \appendix
% \section{Additional Rate-Distortion Detail}
% \prompt{Optional: move SST residual-coding-variant detail here if
% Section~\ref{sec:rd-comparisons} got too long.}

\end{document}
